\subsection{All odd multipair terminals at a root (1,2,v)}\label{proof:Q5-C177}
% Common local aliases for the former repeated C168.1--C168.3 prefaces.
Use the original-measure calculus of
Section~\ref{sec:odd-common-calculus} and the complete root-pair
certificates of Section~\ref{sec:odd-common:fourth}.
For a root \((1,2,v)\), write \(m=v-2\), \(G_m=E+zD\),
\(\beta=\beta_{v-1}=\int D\,d\nu\),
\(\beta_v=\int E\,d\nu\),
\(P_+=h_2f_v\), \(P_-=\bar f_2W_v\),
\(\ell_+=\ell_a\), and \(\ell_-=\ell_s\).
The two vectors, their signed means and full constants are
\eqref{eq:odd-common:root-standard}--\eqref{eq:odd-common:root-losses}.
At \(v=3\), \((E,D)=(1,0)\) and \(\beta=0\).
All subsequent direction choices and packet estimates retain their
stated integer and parameter ranges.
The three complete finite base-vector proofs used below are C173.6.1--C173.6.3 of Section~\ref{proof:Q5-C173}; the negative-root joint payment is C169.5 of Section~\ref{proof:Q5-C169}. The ordinary-fourth outer bridge is Q5-C084, Section 2, with cost $11d$ for every $M\ge1$. The high-domain regrouping uses exactly the fixed-direction terminal triples proved in Sections~\ref{proof:Q5-C169}, \ref{proof:Q5-C172} and \ref{proof:Q5-C173}.
\paragraph{C177.0 Fixed original data and direction}

The construction retains all original arms, equal-arm multiplicities, subdivisions, independent positive core lengths, edge counts and arbitrary original label laws.

% Replace a pure model-definition block only; keep local root/terminal data.
Use the original finite-space model and arm messages of
Section~\ref{sec:odd-common-calculus}.
When the original \(K\)-notation is used, \(K=-B\) and
\(g=K\one=-u\); in either notation
\(\delta_j=\|B^j\one\|_2^2=\|K^j\one\|_2^2\).
All means, norms, adjoints and spectral integrals use the same
original probability measure.
The root residual is (1, 2, v), v at least three; the terminal residual is
\[
(1,a_1,b_1,\ldots,a_j,b_j),\quad
1<a_1<b_1<\cdots<a_j<b_j,\quad j\ge2.
\]
Original positive L, M are independent.

Set the terminal direction as follows:
\[
F=\begin{cases}
qf_2h_{b_2}h_3f_4\prod_{i=3}^jh_{a_i}f_{b_i},
 &(a_1,b_1,a_2)=(2,3,4),\\
qf_{a_1}h_{b_1}\prod_{i=2}^jh_{a_i}f_{b_i},&\text{otherwise}.
\end{cases}                                                    \tag{C177.T184.1}
\]
An empty product is one. The same crossed rule is used at \(j=2\) and every larger pair count. The actual root unit arm gives \(p\); fix
\[
P=\begin{cases}\bar f_2W_v,&a_1=2,\\
h_2f_v,&a_1\ge3,\ L=1,\ v\ge4,\\
\bar f_2W_v,&a_1\ge3\text{; otherwise},\end{cases}
\qquad H=pP\,T^LF.                                          \tag{C177.T184.2}
\]
Inner edges point terminal to root, and outer edges root to receiver. The four arm types retain their specified literal directions. 

For every listed integer, legal finite original host and genuine equal-pair decoration,
\[
\boxed{\begin{gathered}
EH^4\le e^{-705d/128},\qquad\\
I_q(T^MH)^3\le e^{-d/64},\quad I_q(Y)=E[q^{-1/3}|Y|^{8/3}].
\end{gathered}}\tag{C177.T184.3}
\]
The rate \(705/128=11/2+1/128\) includes the actual unit root arm and complete pair.


The original finite calculus gives positive \(A,\allowbreak \allowbreak  T\) with kernels in \([0,\allowbreak \allowbreak  2]\), \(0\le d\le1/3\), \(|u|\le1\), \(\|B\|^2\le\theta\), \(|BV(x)|\le\|V\|\), \(d^2\le \delta_2\), and \(\delta_{r+1}\le\theta \delta_r\). Consistently directed path means are at most one, and for \(n\ge3\) at most \(1-2d\). Genuine equal pairs lie in \([0,\allowbreak \allowbreak  1]\). Delete zero masses only; then \(p,\allowbreak \allowbreak  q\ge\mu(x)>0\), and constant labels give strict positivity.

The common full-vector exponential and all-length ordinary fourth interfaces are
\[
\begin{gathered}
Ee^{4BV}\le e^{4E(BV)+(17/4)\|V\|^2}\quad(V\text{any original real vector}),\\
\qquad I_q(T^MY)^3\le e^{11d}(E|Y|^4)^2\quad(M\ge1).
\end{gathered}\tag{C177.T184.4}
\]
Use the signed mean.

For each actual pair with \(4\le s<t,\allowbreak \) the Q5-C096 envelope is
\[
\begin{gathered}
h_sf_t\le\exp\{R_{st}-16d/5\},\qquad\\
R_{st}=\sum_{r=s+1}^tK^r1.
\end{gathered}\tag{C177.T184.5}
\]
Its \(16/5\) loss is the finite unequal-pair payment.

\paragraph{C177.2 One finite zero-one mask for every later pair}

Put \(G_n(z)=\sum_{r=0}^{n-1}z^r\), \(\beta_1=\beta_2=0\), and for \(n\ge3\),
\[
\begin{gathered}
\beta_n=\sum_{r=1}^{\lfloor(n-1)/2\rfloor}(-1)^{r-1}\delta_r,\\
\qquad\tau=d-\delta_2\le\beta_n\le d.
\end{gathered}
\]
The exact finite convolution is
\[
1+G_n(B)u=\bar f_n+\sum_{r=3}^n\beta_r \bar f_{n-r}.                \tag{C177.T184.6}
\]
The spectral measure has mass d, support [0, theta], and moments integral \(x^r=\delta_{r+1}.\) At \(z=i\) sqrt(x),
\[
\begin{gathered}
\|(A_0+zC_0)(B)u\|^2=\int(A_0^2+xC_0^2),\\
\quad E[(A_0+zC_0)(B)u]=-\int C_0.
\end{gathered}\tag{C177.T184.7}
\]
For \(n\ge2\), both coordinates of \(G_n\) lie in \([1-x,\allowbreak \allowbreak  1]\), while \(n=1\) gives \((1,\allowbreak \allowbreak  0)\). For \(n\ge3\), the real coordinate is at most \(1-x+x^2\) and the modulus at most one. Retain \(|G_2|\le\sqrt{1+\theta}<6/5\).

For any finite group of later pairs, their exponent intervals are disjoint; their sum has form
\[
S=-P(-B)u,\qquad P(z)=\sum_{i}\sum_{k=s_i}^{t_i-1}z^k
\]
Thus P is a finite zero-one mask, with minimum shorter length s. Four residue classes give
\[
\begin{gathered}
|P(i\sqrt x)|^2\le\frac{1+x}{(1-x^2)^2}\\
\le\frac{32768}{16245}<\frac{101}{50}<\frac{100}{49}.
\end{gathered}\tag{C177.T184.8}
\]
After removing the lowest \(z^s\), its remaining mask has norm less than \(10/7\), independently of pair count. Row and spectral bounds give
\[
\begin{gathered}
|S|\le(10/7)\theta^{(s-1)/2}\sqrt d,\\
\quad \|S\|\le(10/7)\theta^{(s-2)/2}\sqrt{\delta_3},\\
\quad |S/u|_{\rm spectral}\le(10/7)x^{s/2}.
\end{gathered}\tag{C177.T184.9}
\]
The final notation is a finite spectral polynomial.

For s at least five, its real and imaginary coordinates satisfy
\[
\begin{gathered}
|A_S|\le\frac{x^3}{1-x^2},\qquad\\
|C_S|\le\frac{x^2}{1-x^2}.
\end{gathered}\tag{C177.T184.10}
\]
For \(s\ge6\), improve the second bound to \(x^3/(1-x^2)\). In each residue class bound the positive and negative subseries separately. The positive original mean terms arise from paired powers \(K^7,\allowbreak \allowbreak  K^8\), then \(K^{11},\allowbreak \allowbreak  K^{12}\), and so on. Hence
\[
\begin{gathered}
ES-E(gS)\le2\sum_{r\ge0}\delta_{4+2r}\\
\le2C_\beta \delta_4<d/6,\quad g=-u,\ C_\beta=1024/855.
\end{gathered}\tag{C177.T184.11}
\]
The geometric tail bounds only the actual finite subseries. We shall also use the exact rational inequalities
\[
\sqrt\theta<16/25<13/20,\quad\theta^2<1/6,\quad
\theta^3<7/100<1/12,
\]
\[
\begin{gathered}
C_\beta\theta=416/855<1/2,\quad\\
C_\beta\theta^2=169/855<1/5,\quad\\
C_\beta\theta^3=2197/27360<1/12,\quad\\
C_\beta\theta^4<1/30.
\end{gathered}\tag{C177.T184.12}
\]
These bounds include zero frequency, every finite highest power and all pair counts.

\paragraph{C177.3 Actual first-row propagation and the full root payment}

If \(F\le qe^{R-\Gamma}\) and \(|R|\le1/2\), use \(e^R\le1+R+3R^2/5\), \(q\le2\) and the original kernel bound two:
\[
\begin{gathered}
TF\le e^{-\Gamma}(\bar f_2+BR+\delta),\\
\quad\delta=ER+E(uR)+\sqrt{dER^2}+(12/5)ER^2.
\end{gathered}\tag{C177.T184.13}
\]
Expand \(T(qR)=ER+E(uR)+BR+B(uR)\), applying Cauchy only to the last row term. If only \(|R|\le3/4\) is available, use \(e^R\le1+R+2R^2/3\), replacing \(12/5\) by \(8/3\). The positive Taylor-tail coefficient is at most \(355/544<2/3\), and the negative integral coefficient at most \(1/2\).

Take \(D\ge\max(\delta,\allowbreak \allowbreak  0)\), \(D_i=EB^iR\), \(\gamma_r=\beta_r-D_{r-2}\). Actual positive finite propagation gives
\[
\begin{gathered}
T^LF\le e^{-\Gamma}Y_L,\\
\quad Y_L=\bar f_{L+1}+B^LR+D\bar f_{L-1}\\
+\sum_{i=1}^{L-1}D_i \bar f_{L-1-i}.
\end{gathered}\tag{C177.T184.14}
\]
After proving every gamma sign, retain the final original \(\bar f_0\) term for \(L\ge2\):
\[
\begin{gathered}
Y_L\le1+u+2D-\gamma_{L+1}+BQ_L(B)u,\\
\quad Q_L=G_L+z^{L-1}r(z),\quad R=r(B)u,\\
\quad E[BQ_L(B)u]=-\gamma_{L+2}.
\end{gathered}\tag{C177.T184.15}
\]
At \(L=1\), use exactly \(Y_1=1+u+D+B(u+R)\). Combine the complete root envelopes of Lemma~\ref{lem:odd-common:root-pairs} with \eqref{eq:odd-common:curvature}, retaining the original $\beta\one$ component.


For L at least two, the complete payment is
\[
\begin{gathered}
\log EH^4\le-4\Gamma+8D-4(\gamma_{L+1}+\gamma_{L+2})\\
-4\ell_-+(17/4)\|[Q_L-zG_n(z)]u\|^2+(10/3)\beta^2.
\end{gathered}\tag{C177.T184.18}
\]
All these actual nonunit rows use the negative root. At \(L=1\), replace \(8D\) minus four times the gamma sum by \(4D-4d+4D_1\), using the prescribed positive or negative root. At \(L=2\), retain \(\gamma_3+\gamma_4=2d-D_1-D_2\). The curvature cost is at most \(5d/96\) on \(d\le1/64\), and \(5d/12\) on \(d\le1/8\), and zero at \(v=3\).

\paragraph{C177.4 Root bases, continuous short rectangles and all finite longer lengths}

Use the full base proofs in C173.6 and C169.5. The negative first-row base has joint cost at most two and modulus below \(8/5\), or \(6/5\) at \(n=1\). The positive first-row base, used for \(n\ge2\), has modulus at most one and joint cost \(119/256<1/2\). The negative second-row base has modulus below \(23/16\), joint cost at most \(3/2\); every longer finite base has the same modulus bound and joint cost at most \(1/4\). These whole-frequency proofs retain \(G_1,\allowbreak \allowbreak  G_2\).


If \(r=-t+zs\) lies in the continuous rectangle
\[
-1/16\le t\le2/3,\qquad-1/5\le s\le1/2,               \tag{C177.T184.19}
\]
the negative-root short rows, including their actual means, satisfy
\[
\begin{gathered}
J_1=4t-4[E_0+(1+x)D_0]\\
+(17/4)[(1-t+xD_0)^2+x(s-E_0)^2]\le4,
\end{gathered}\tag{C177.T184.20}
\]
\[
\begin{gathered}
J_2=4t+4xs-4[E_0+(1+x)D_0]\\
+(17/4)[(1+x(D_0-s))^2+x(1-t-E_0)^2]\le6.
\end{gathered}\tag{C177.T184.21}
\]
The \(D_0\)-derivative of \(J_1\) is at most \(-4(1+x)+(17/2)x(17/16+x)\), whose endpoint is \(-1133/2048<0\). The square is dominated by \(s=-1/5\). Convexity in \(t\) reduces to endpoints; the difference bracket between \(2/3\) and \(-1/16\) is at most \(-371/192\), so the negative endpoint is larger. For \(n\ge2\), set \(D_0=1-x\), use \(1-x\le E_0\le1\) and \(x-x^2\le1/4\), giving \(J_1\le98373/25600<4\). At \(n=1\), direct substitution gives \(77673/25600<4\). This covers the whole continuous rectangle.

The expression \(J_2\) decreases in \(D_0,\allowbreak \allowbreak  E_0,\allowbreak \allowbreak  s\); the first derivative has endpoint upper bound \(-2501/5120<0\). For \(n\ge2\), set both root coordinates to \(1-x\), and \(s=-1/5\). Convexity in \(t\) selects \(2/3\), since its difference bracket is \(4+(17/4)x[29/48+2(E_0-1)]>0\). The real coordinate is at most \(4/3\), using \(\theta(6/5-\theta)=1651/5120<1/3\), and the imaginary coefficient has absolute value at most \(2/3\). Thus \(J_2\le19007/3840<6\). At \(n=1\), the square is below \(7/5\) and scalar at most \(-4/3\), giving \(J_2<277/60<6\). If \(0\le t,\allowbreak \allowbreak  s\le2/3\), enlarge to \([0,\allowbreak \allowbreak  1]^2\): \(s=0\) dominates, the first-row \(t=0\) and second-row \(t=1\) endpoints give \(J_1\le2,\allowbreak \allowbreak  J_2\le6\).

For every longer \(L\), a modulus bound \(\rho\) on \(r\) bounds the perturbation by \(\rho x\); its full added squared-norm cost is
\[
\mathcal E(\rho)d=(17/4)[(23/8)\rho\theta+\rho^2\theta^2]d. \tag{C177.T184.22}
\]
For \(\rho=1/2,\allowbreak \allowbreak  2/3,\allowbreak \allowbreak  3/5\), use respectively \(8/3,\allowbreak \allowbreak  11/3,\allowbreak \allowbreak  13/4\). This controls every finite longer length.

\paragraph{C177.5 Case I: minimum first pair at least three}

Take the first canonical pair \(qf_{a_1}h_{b_1}\le qe^{R_0-\Gamma_0}\). At \(a_1=3\), \(\Gamma_0=d+\beta_{b_1}+\beta_{b_1-1}\ge35d/16\), \(R_0=-\beta_{b_1-1}u-B^3G_{b_1-3}(-B)u\), with mixed mean nonpositive. For \(b_1=4,\allowbreak \allowbreak  5\), the means are \(-\delta_2-d^2\), \(-\delta_2-d^2+\delta_3\); for \(b_1\ge6\), at most \(-3\delta_2/16-\kappa d^2\). For \(a_1\ge4\), the full four-beta envelope is
\[
\begin{gathered}
\Gamma_0=\beta_{a_1}+\beta_{a_1-1}+\beta_{b_1}+\beta_{b_1-1}\ge19d/8,\\
\quad r_0=-c+(-1)^{a_1}z^{a_1}G_{b_1-a_1}(-z),\\
\ c=\beta_{b_1-1}-\beta_{a_1-1},\ |c|\le \delta_2.
\end{gathered}
\]
This comes from the actual finite convolution retaining \(f_0,\allowbreak \allowbreak  f_1\) and the \(q\) term. Its mixed mean is at most \(\delta_2d+2\delta_3\le7d/20\) on the present low domain.

Let S be all later pair intervals, starting at \(s=a2\) at least five. The actual pair losses give
\[
\begin{gathered}
F\le q e^{R-\Gamma},\quad R=R_0+S,\quad\\
\Gamma\ge(35/16+16/5)d=(431/80)d.
\end{gathered}\tag{C177.T184.23}
\]
Each additional pair increases the scalar loss while the same single mask bounds \(S\). At \(a_1=3\), \(|r_0|\le1/64+507/1600\), while the later mask modulus is below \(13/80\), giving \(|r|<99/200<1/2\). At \(a_1\ge4\), \(|r_0|\le\theta/64+(6/5)\theta^2\), and the later mask starts at six, again giving modulus below \(1/2\). The original row bound gives \(|R|\le d+2\sqrt d\le17/64<1/2\). The mixed mean is at most \(7d/20+d/6=31d/60\), and \(ER^2\le d/4\). Therefore
\[
\delta/d\le31/60+1/2+3/5=97/60<5/3.
\]
Take \(D=5d/3\). Spectral Cauchy gives \(|D_i|\le\rho\theta^{(i-1)/2}d\); hence \(\gamma_3\ge d/2\), \(\gamma_4\ge27d/40\), and later gammas at least \(25d/64>d/3\).

At \(L=1\), only \(v=3\) uses the negative root. Base cost at most two, modulus at most \(6/5\), and \(\rho=1/2\) add norm cost \(493d/80\); including \(4D_1\le2d\), the total is below \(21d/2\). At \(L=2\), perturbation modulus is below \(13/40\), added cost below \(9d/2\); with base \(3d/2\), mean cost below \(33d/10\), the total is at most \(93d/10\). Longer lengths add \(8d/3\).

\paragraph{C177.6 Case II: first pair (2, b1), b1 at least four}

The actual first-pair convolution gives
\[
\begin{gathered}
\Gamma_0=\beta_{b_1}+B_0\ge2\kappa d,\\
\quad R_0=B_0g-\sum_{r=3}^{b_1}K^r1,\\
\quad B_0=\beta_{b_1-1}.
\end{gathered}
\]
The later mask starts at \(a_2\ge5\), and \(\Gamma\ge351d/80\). The mixed mean of \(R_0\) is at most \(-2\delta_2\) at \(b_1=4\), and \(-2\kappa \delta_2\) otherwise. Retaining the finite \(G_2\) exception gives, with \(v_0=\delta_2/d\),
\[
\begin{gathered}
\sqrt{ER^2/d}\le d+C\sqrt{\theta v_0},\qquad\\
C=8/5\ (b1=4),\quad C=5/4\ (b1\ge5).
\end{gathered}\tag{C177.T184.24}
\]
For the first coefficient use \(\sqrt{1+\theta}+(10/7)\theta^{3/2}<6/5+169/448<8/5\). For the second use \(1+(10/7)\theta^2<5/4\), since \(a_2\ge6\). The positive mixed-mean tail of \(S\) is at most \(2C_\beta\theta^2\delta_2\), so
\[
\delta/d\le-\lambda v_0+d+C\sqrt{\theta v_0}
+(12/5)[d+C\sqrt{\theta v_0}]^2,
\]
Take \(\lambda_4=1372/855\), \(\lambda_5=2\kappa-338/855\). In the variable \(\sqrt{v_0}\), the quadratic and linear coefficients are positive, so the bound increases on \([0,\allowbreak \allowbreak  \theta]\); it also increases in \(d\). At \(b_1=4\), the endpoint square root is \(213/320\), and the value \(221307/128000-4459/6840<11/10\), with cross-product margin \(20084120\). At \(b_1\ge5\), the square root is \(67/128<21/40\); use \(\lambda_5>3/4,\allowbreak \allowbreak  \theta>2/5\) to get \(2373/2000-3/10=1773/2000<9/10\). Uniformly take \(D=11d/10\).

The first-pair polynomial is \(r_0=-B_0-xU+zxV\), with \(G_{b_1-2}(-z)=U-zV\). The later mask gives \(|r|<1/64+39/80+13/80=213/320<2/3\), and \(|R|\le d+2\sqrt d<1/2\). Its coordinates satisfy
\[
0\le t\le2/3,\qquad0\le s\le2/3.                     \tag{C177.T184.25}
\]
The lower bounds follow from \(x(1-x)-C_\beta x^3>0\), \(x(1-x)-C_\beta x^2>0\); upper bounds from \(1/64+C_\beta\theta<2/3\), \(\theta+C_\beta\theta^2<2/3\). Thus the full short-window bounds apply.

The finite moments obey \(|D_1|\le d^2+\delta_2+C_\beta \delta_4<2d/3\), \(D_2\le \delta_3+C_\beta \delta_4<d/4\), and every later positive part is at most \(B_0\delta_2+\delta_3+C_\beta \delta_5<d/4\). Use \(\theta/64<1/128,\allowbreak \allowbreak  \theta^2<1/6,\allowbreak \allowbreak  C_\beta\theta^4<1/30\). Hence \(\gamma_3\ge d/3,\allowbreak \allowbreak  \gamma_4\ge3d/4\), and later gammas are at least \(11d/32>d/3\). The long cost is \(11d/3\).

\paragraph{C177.7 Case III: first pair (2, 3), next shorter length at least five}

The actual terminal is \(qf_2h_3\prod_{i\ge2}h_{a_i}f_{b_i}\). The exact first-pair sum is \(f_2+h_3=2-d-K^3\one\), so \(\Gamma_0=d,\allowbreak \allowbreak  R_0=-K^3\one\). Later pairs give
\[
\Gamma\ge21d/5,\quad R=-K^31+S,\quad S\text{ first shorter length }s=a2\ge5.
\]
The mixed mean is at most \(-(1-2C_\beta\theta^2)\delta_2\). For \(v_0=\delta_2/d\),
\[
ER^2/d\le\theta(617/448)^2v_0,
\]
Here \(1+(10/7)\theta^{3/2}<617/448\). The first-row bound increases in \(\sqrt{v_0}\), with positive quadratic and linear coefficients. At \(v_0=\theta\), the square root is below \(9/16\), and the mean deduction exceeds \(6d/25\). Thus
\[
\delta/d<9/16+243/320-6/25=8655/8000<11/10.
\]
Take \(D=11d/10\), \(\rho=3/5\). The original pointwise row bound is still below \(1/2\).

The finite coordinates \(r=-x+A_S+zC_S\) give \(0\le t\le416/855<2/3\), \(|s|\le169/855<1/5\). Also \(D_1\le \delta_2+C_\beta \delta_4<d/2\), \(D_2\le C_\beta \delta_4<d/12\). Across the four residue classes, later positive parts start at \(C_\beta \delta_5,\allowbreak \allowbreak  C_\beta \delta_5,\allowbreak \allowbreak  \delta_4+C_\beta \delta_6,\allowbreak \allowbreak  C_\beta \delta_6\), then decrease by \(\theta^2\) every four indices. They are at most \(\delta_4+C_\beta \delta_5<d/10\), using \(\theta^3<43/640,\allowbreak \allowbreak  C_\beta\theta^4<21/640\). Thus \(\gamma_3\ge d/2,\allowbreak \allowbreak  \gamma_4\ge11d/12\), and later gammas at least \(79d/160\ge15d/32\). Use short costs four and six; the long adjacent-gamma sum is at least \(15d/16\), with norm increment \(13d/4\).

\paragraph{C177.8 Cases IV and V: one fixed crossed direction}

Let \(e=b_2\ge5\). The first two pairs are always \(qf_2h_eh_3f_4\); later pairs retain \(h_{a_i}f_{b_i}\). The retained coefficients are those of Lemma~\ref{lem:unit-retained-pair} with $\sigma=9/4$; Lemmas~\ref{lem:reference-retained:strong-norm} and \ref{lem:reference-retained-packet} give their actual continuous frequency and error bounds, here on the smaller domain \(d\le1/8\). At \(e=5\), take \(c_1=d,\allowbreak \allowbreak  c_2=d/4+\delta_2,\allowbreak \allowbreak  c_3=0\); otherwise use


\[
\begin{gathered}
c_1=\beta_{e-1},\quad c_2=\lambda\beta_{e-2},\quad c_3=\lambda\beta_{e-3},\\
\quad\lambda=\frac{9d/4-\beta_e-\beta_{e-1}}{\beta_{e-2}+(1-d)\beta_{e-3}}\in[0,1].
\end{gathered}
\]
Each retained coefficient lies within the actual positive convolution, with positive denominator and total available constant exceeding \(9d/4\). Put \(C=c_1+c_2+c_3,\allowbreak \allowbreak  D_*=c_2+c_3,\allowbreak \allowbreak  E=c_3\). Then \(f_2+h_e\le2-9d/4-U_e\), where \(U_e=(A_{\rm old}-zB_{\rm old})u\), \(A_{\rm old}\le3/8+9d/5\). For \(e\ge6\), its coordinates are \(C+x(a_f-E),\allowbreak \allowbreak  D_*+xb_f\); for \(e=5\), they are \(5d/4+\delta_2+x-x^2,\allowbreak \allowbreak  d/4+\delta_2+x\). Retain \(C\ge5d/4+dE,\allowbreak \allowbreak  D_*\ge d/4,\allowbreak \allowbreak  D_*\le(17/16+d)d,\allowbreak \allowbreak  0\le E\le d\).

Combine this with the exact \(h_3+f_4=2-3d+K^4\one-dg\), yielding
\[
\begin{gathered}
F_{\rm cross}\le qe^{R_c-21d/4},\quad\\
R_c=(-A_c+zB_c)u,\quad A_c=A_{old}-d,\quad B_c=B_{old}-x.
\end{gathered}
\]
The second coordinate may be negative. Its complete range is
\[
\begin{gathered}
0\le A_c\le19/40,\quad -x^2\le B_c\le19d/16\le19/128,\\
\quad A_c+B_c\ge d/2+x/16.
\end{gathered}\tag{C177.T184.26}
\]
The final inequality uses \(a_f+b_f\ge2(1-x)\), \(1-2x-E\ge1/16\) for \(e\ge6\); substitute directly at five. Hence \(ER_c+E(uR_c)=-\int(A_c+B_c)\le0\), and \(\rho_c^2<(19/40)^2+\theta/36<1/4\). The original row formula is \(R_c=(C-d)g+D_*Kg+(E-1)K^2g-\sum_{r=5}^eK^r\one\), so
\[
|R_c|\le25/256+(19/128+119/125)\sqrt{1/8}<1/2.
\]
Using \(\sqrt{1/8}<9/25\), its second term is below \(158463/400000\), still giving total below \(1/2\). Thus \( \delta_c<d/2+3d/5=11d/10\). The moments obey \(D_1\le19d/40,\allowbreak \allowbreak  D_2<d/16\), and later positive parts below \(d/12\): \(D_3\le0,\allowbreak \allowbreak  D_4\le \delta_5,\allowbreak \allowbreak  D_5<19\delta_3/40,\allowbreak \allowbreak  D_6<19d\delta_4/16\), with each residue class then decreasing. Hence \(\gamma_3\ge d/2,\allowbreak \allowbreak  \gamma_4\ge15d/16\), and later gammas at least \(49d/96\).

In IV, \(j=2\), use \(\Gamma=21d/4,\allowbreak \allowbreak  D=11d/10,\allowbreak \allowbreak  \rho=1/2\). The coordinates satisfy the full short rectangle, giving costs four and six, with long increment \(8d/3\).

In V, \(j\ge3\), the later-pair mask starts at \(a_3\ge6\):
\[
R=R_c+S,\qquad\Gamma\ge(21/4+16/5)d=169d/20.
\]
The later mask has \(|S|<13/224<1/12\), so \(|R|<7/12<3/4\). Use the Taylor coefficient \(2/3\). Its spectral modulus is less than \(1/2+(10/7)\theta^3<3/5\), and mixed mean at most \(d/6\). Therefore
\[
\delta/d\le1/6+3/5+(8/3)(9/25)=259/150<7/4.
\]
Take \(D=7d/4\). The minimum-six coordinate bounds give \(D_1<19d/40+d/12<2d/3\), \(D_2<d/16+d/30<d/3\), and later positive parts below \(d/12+d/30<d/8\). Hence \(\gamma_3\ge d/3,\allowbreak \allowbreak  \gamma_4\ge2d/3\), and later gammas at least \(15d/32\); the long adjacent sum remains at least \(15d/16\).

For V at the short rows use the complete coarse joint payment. At \(L=1\), base cost at most two and modulus below \(8/5\), with \(\rho=3/5\), add at most \(969d/100\); including \(4D_1<8d/3\), the total is below \(29d/2\). At \(L=2\), additional norm cost is below \(11d/2\), base \(3d/2\), and mean cost below \(4d\), giving total at most \(11d\). Longer lengths add \(13d/4\). One finite mask controls every later pair count.

\paragraph{C177.9 Complete rational payments on every low domain}

The table lists positive fourth losses. Cases I--II use \(d\le1/64\); III--V use \(d\le1/8\).


\begingroup\small
\setlength{\tabcolsep}{3pt}
\setlength{\LTleft}{0pt}\setlength{\LTright}{\fill}
\begin{longtable}{@{}p{0.2200\linewidth}p{0.7000\linewidth}@{}}
\caption{All odd multipair terminals at a root (1,2,v): parameter values.}\label{tab:parameters:q5-c177:1}\\
\toprule
Entry & Lower Gamma/d \\
\midrule\endfirsthead
\toprule
Entry & Lower Gamma/d \\
\midrule\endhead
\bottomrule\endfoot
\(I,\allowbreak a_1\ge 3\) & \(431/80\); D/d: \(5/3\); Curvature cost/d: \(5/96\); \(L=1\) loss: \(1333/160\); \(L=2\) loss: \(659/96\); Every \(L\ge 3\) loss: \(3799/480\). \\
\addlinespace[2pt]
\(II,\allowbreak a_1=2,\allowbreak b_1\ge 4\) & \(351/80\); D/d: \(11/10\); Curvature cost/d: \(5/96\); \(L=1\) loss: \(6287/480\); \(L=2\) loss: \(1027/96\); Every \(L\ge 3\) loss: \(715/96\). \\
\addlinespace[2pt]
\(III,\allowbreak a_1=2,\allowbreak b_1=3,\allowbreak a_2\ge 5\) & \(21/5\); D/d: \(11/10\); Curvature cost/d: \(5/12\); \(L=1\) loss: \(719/60\); \(L=2\) loss: \(115/12\); Every \(L\ge 3\) loss: \(47/6\). \\
\addlinespace[2pt]
IV, cross, \(j=2\) & \(21/4\); D/d: \(11/10\); Curvature cost/d: \(5/12\); \(L=1\) loss: \(971/60\); \(L=2\) loss: \(827/60\); Every \(L\ge 3\) loss: \(757/60\). \\
\addlinespace[2pt]
V, cross, \(j\ge 3\) & \(169/20\); D/d: \(7/4\); Curvature cost/d: \(5/12\); \(L=1\) loss: \(953/60\); \(L=2\) loss: \(983/60\); Every \(L\ge 3\) loss: \(589/30\). \\
\end{longtable}
\endgroup


The common full-vector payment gives every entry by the following rational arithmetic.
\[
431/20-20/3+4-21/2-5/96=1333/160,
\]
\[
431/20-40/3+8-93/10-5/96=659/96,
\]
\[
431/20-40/3+8/3-1/4-8/3-5/96=3799/480.
\]
For II:
\[
351/20-22/5+4-4-5/96=6287/480,
\]
\[
351/20-44/5+8-6-5/96=1027/96,
\]
\[
351/20-44/5+8/3-1/4-11/3-5/96=715/96.
\]
For III:
\[
84/5-22/5+4-4-5/12=719/60,
\quad84/5-44/5+8-6-5/12=115/12,
\]
\[
84/5-44/5+15/4-1/4-13/4-5/12=47/6.
\]
For IV:
\[
21-22/5+4-4-5/12=971/60,
\quad21-44/5+8-6-5/12=827/60,
\]
\[
21-44/5+15/4-1/4-8/3-5/12=757/60.
\]
For V:
\[
169/5-7+4-29/2-5/12=953/60,
\quad169/5-14+8-11-5/12=983/60,
\]
\[
169/5-14+15/4-1/4-13/4-5/12=589/30.
\]
The weakest low-domain loss is \(659/96\), exceeding \(705/128\) by \(521/384\). Every bound is a continuous-frequency finite-mask estimate.

\paragraph{C177.10 Same-direction regrouping on the high domains}

For every actual pair with s at least four, the original row bound gives
\[
\begin{gathered}
h_sf_t\le\exp\{(39/125)\sqrt d-16d/5\}\\
\le e^{-88d/125}\le1\quad(d\ge1/64).
\end{gathered}\tag{C177.T184.27}
\]
Its tail is at most \((6/5)\theta^{(s-1)/2}\sqrt d<39\sqrt d/125\), and \(\sqrt d\le8d\) when \(d\ge1/64\).

For I and II on \(d\ge1/64\), drop only the paid later pair factors from the upper bound, obtaining \(H\le pPT^L(qf_{a_1}h_{b_1})\). Case I uses C172 or C173 according as \(a_1=3\) or at least four; case II uses C169 with its fixed negative root. Every rate is \(705/128\).

For III on d at least one eighth, retain every arm direction and regroup the same factors as
\[
F=(qf_2h_{a_2})(h_3f_{b_2})\prod_{i=3}^jh_{a_i}f_{b_i}. \tag{C177.T184.28}
\]
This only changes the multiplication order at one root. Here \(b_2\ge6\), and the actual centered-row estimate below gives \(h_3f_{b_2}\le e^{-d}\); other later pair factors are at most one. Compare with \(qf_2h_{a_2}\), \(a_2\ge5\), with the same negative root, and use C169.

For IV and V on the same high domain, the original product is already
\[
F=(qf_2h_{b_2})(h_3f_4)\prod_{i=3}^jh_{a_i}f_{b_i}.     \tag{C177.T184.29}
\]
The factor \(h_3f_4\le e^{-d}\), and later pairs start at six, hence are at most one. Compare with \(qf_2h_{b_2}\), \(b_2\ge5\), and use the same C169 negative-root row. The crossed orientation is identical in IV and V.

The actual centered row is as follows. At \(c\ge5\), put \(B_0=\beta_{c-1}\), \(\Gamma_c=d+\beta_c+B_0\), \(U=\sum_{r=3}^{c-1}K^r\one\), \(a=B_0-EU=\Gamma_c-2d\ge0\). The arm sum gives \(h_3f_c\le e^{-\Gamma_c-R_0}\), \(R_0=ag-K(U-EU)\). Centered-row Cauchy and \(g^2+pq=1\) give \(|R_0|\le\sqrt{a^2+\operatorname{\operatorname{Var}}U}\). For \(c\ge6\), \(\operatorname{\operatorname{Var}}U\le\theta^2d\), \(a\ge3d/16\); when \(d\ge1/8\), this variance is less than \(11d^2/8\le d^2+2ad\). At \(c=5\), retain \(\operatorname{\operatorname{Var}}U\le\theta^2(1+\theta)d\), \(a\ge\kappa d\), and \(8\theta^2(1+\theta)=7605/4096<35/16\le1+2a/d\). Thus \(|R_0|\le a+d\), leaving loss at least \(d\). At four, use \(h_3+f_4=2-3d+K^4\one-dg\); the centered row bounds its remainder by \(\sqrt{d^2+\delta_3}\le2d\) on the same domain. This proves \(h_3f_4\le e^{-d}\) as well.

The low and high intervals share their closed endpoints, with the same fixed directions. At zero, original leaf integration gives equality, allowing nonzero kernel. Apply the Q5-C084 all-length ordinary fourth bridge:
\[
I_q(T^MH)^3\le e^{11d}(EH^4)^2
\le\exp\{[11-2(705/128)]d\}=e^{-d/64}.
\]
